% !PRINTED 268
\weberpage{268}
Accordingly (11) becomes
\[
(uv)^{\frac{p+1}{2}}\left[\Bigl(\frac{v}{u}\Bigr)^{r\alpha}+\Bigl(\frac{u}{v}\Bigr)^{r\alpha}\right]\left[(uv)^{s\beta}+\frac{\varepsilon^\beta2^{s\beta}}{(uv)^{s\beta}}\right],
\]
and (12)
\[
(uv)^{\frac{p+1}{2}}\left[(uv)^{s\beta}+\frac{\varepsilon^\beta2^{s\beta}}{(uv)^{s\beta}}\right].
\]
Now the well-known formulae hold, when $x, \gamma$ are arbitrary
quantities and
\[
x+\frac{\gamma}{x}=y
\]
is put,
\[
x^2+\frac{\gamma^2}{x^2}=y^2-2\gamma, \qquad x^3+\frac{\gamma^3}{x^3}=y^3-3\gamma y,\ \ldots,
\]
\[
\Bigl(x^n+\frac{\gamma^n}{x^n}\Bigr)\Bigl(x+\frac{\gamma}{x}\Bigr)=x^{n+1}+\frac{\gamma^{n+1}}{x^{n+1}}+\gamma\Bigl(x^{n-1}+\frac{\gamma^{n-1}}{x^{n-1}}\Bigr),
\]
from which one recognises, by the inference from $n$ to $n+1$, that
\[
x^n+\frac{\gamma^n}{x^n}
\]
can be represented, for every arbitrary $n$, as a polynomial of degree
$n$ in $y$ which, when $\gamma$ is an integer, has integer coefficients,
its leading coefficient being $=1$.

The two quantities
\[
\Bigl(\frac{v}{u}\Bigr)^{r\alpha}+\Bigl(\frac{u}{v}\Bigr)^{r\alpha} \qquad \text{and} \qquad (uv)^{s\beta}+\frac{\varepsilon^\beta2^{s\beta}}{(uv)^{s\beta}}
\]
can therefore be represented in this way as polynomials in $A$ and $B$
of degrees $\alpha$ and $\beta$, and we find, if we
further separate off the term corresponding to the value
$\beta=\dfrac{p-1}{2s}$, which has the coefficient $-1$, and denote
integer coefficients by $c_{\alpha,\beta}$:
\weq{14}{\frac{\Phi_p(u,v)}{(uv)^{\frac{p+1}{2}}}=A^{\frac{p+1}{2r}}-B^{\frac{p-1}{2s}}+\sum^{\alpha,\beta}c_{\alpha,\beta}A^\alpha B^\beta,}
in which $\alpha$ and $\beta$ are bound to the limits
\[
0\leqq\alpha<\frac{p+1}{2r}, \qquad 0\leqq\beta<\frac{p-1}{2s}
\]
This is the form which we gave to the modular equations up to $p=19$
in the preceding section.
